\bigskip
\bigskip
\subsubsection*{Խնդիրներ և հարցեր թեմա 15-ի վերաբերյալ}

\begin{enumerate}[label=\thesection.\arabic*.]

\red{Evkilyan te Evkidlesay? Who knows? Or more importantly` who cares? Definitely not me.}'}

% 15.1
\item Ճի՞շտ է պնդումը. Կապակցված տարածության ցանկացած ենթատարածություն կապակցված է:

% 15.2
\item Ապացուցեք. Կապակցված $X$ տարածության ցանկացած $X/R$ ֆակտոր-տարածություն կապակցված է:

% 15.3
\item Ապացուցեք. $(\mathbb{R}, \text{սովոր.})$ տարածության ռացիոնալ կետերի ամեն մի ենթատարածություն կապակցված չէ:

% 15.4
\item Ունենք $f: X \to Y$ անընդհատ արտապատկերում: Ճի՞շտ է պնդումը. $Y$-ում կապակցված ամեն մի $B$ ենթաբազմության $f^{-1}(B)$ նախապատկերը կապակցված է $X$-ում:

% 15.5
\item Ապացուցեք. $\mathbb{R}^n$ էվկլիդեսյան տարածության ամեն մի ուռուցիկ ենթաբազմություն կապակցված ենթատարածություն է:
\begin{hint}
Սևեռելով $X \subset \mathbb{R}^n$ ուռուցիկ ենթաբազմության որևէ $x_0$ կետ դիտարկեք $x_0$ սկզբնակետով բոլոր հատվածները, որոնք ընկած են $X$-ում:
\end{hint}

% 15.6
\item $\mathbb{R}^3$ էվկլիդեսյան տարածության հետևյալ երեք $\{(x,y,z); x^2+y^2-z^2=\delta, \text{որտեղ } \delta=-1, 0, 1\}$ մակերևույթներից որո՞նք են կապակցված (պատասխանը հիմնավորեք):

% 15.7
\item Ճի՞շտ է արդյոք պնդումը. $(\mathbb{R}, \rightarrow)$ տարածության հետևյալ ենթաբազմություններից՝ ա) $\{0,1\}$, բ) $[0,1)$, գ) $[0,1]$, դ) $(0,0)$, ե) $(0,1]$ ոչ մեկը կապակցված չէ:

% 15.8
\item Ապացուցեք. $\mathbb{R}^2$ էվկլիդեսյան հարթության այն բոլոր կետերի $A$ ենթաբազմությունը, որոնց կոորդինատներից գոնե մեկը ռացիոնալ թիվ է, կապակցված ենթաբազմություն է:
\begin{hint}
Հիմնավորեք, որ $A$-ի ցանկացած երկու կետեր կարող են միացվել իրար երկու կամ երեք հատվածներից կազմված բեկյալով, որն ամբողջությամբ ընկած է $A$-ի մեջ:
\end{hint}

% 15.9
\item Ճի՞շտ է արդյոք, որ $\mathbb{R}^2$ հարթության այն բոլոր կետերի $B$ ենթաբազմությունը, որոնց կոորդինատներից միայն մեկն է ռացիոնալ թիվ, կապակցված ենթաբազմություն է:
\begin{hint}
Ըստ պայմանի $B$-ն չունի ընդհանուր կետ $y=x$ ուղղի հետ: Ուստի $B$-ն ներկայացվում է որպես երկու ենթաբազմությունների միավորում, որոնցից մեկը ընկած է այդ ուղղից վերև, իսկ մյուսը՝ ներքև:
\end{hint}

% 15.10
\item Ապացուցեք. $\mathbb{R}$ թվային ուղղի որևէ $[a,b]$ ենթատարածություն հոմեոմորֆ չէ $\mathbb{R}^2$ հարթության որևէ $[c,d] \times [e,f]$ ենթատարածության:

% 15.11
\item Ապացուցեք. $X$ տոպոլոգիական տարածությունը կապակցված է այն և միայն այն դեպքում, երբ ցանկացած $f: X \to Y$ անընդհատ արտապատկերում, որտեղ $Y$-ը մեկից ավելի կետեր պարունակող դիսկրետ տարածություն է, հաստատուն արտապատկերում է:

% 15.12
\item Ապացուցեք. $(X, \tau)$ տոպոլոգիական տարածությունը կապակցված չէ այն և միայն այն դեպքում, երբ գոյություն ունի $f: (X, \tau) \to (Y, \text{դիսկր.})$ անընդհատ սյուրյեկտիվ արտապատկերում, որտեղ $Y$-ը երկկետանոց բազմություն է:

% 15.13
\item Ապացուցեք. եթե $A$-ն $X$ տարածության կապակցված ենթատարածություն է և $A \subset Y \subset \bar{A}$, ապա $Y$-ը $X$-ի կապակցված ենթատարածություն է:

\end{enumerate}
